\begin{center}
\textsc{Gauss} an \textsc{Olbers}.\quad Göttingen, 14.~Juni 1830.
\end{center}

\par\noindent\hspace*{\parindent}\mbox{.\ .\ .\ .\ .\ .}\quad Die Zeit, die mir in diesem Sommer für eigene Arbeit bleibt, will ich der Fortsetzung meiner Abhandlung über die biquadratischen Reste widmen, damit diese Arbeit, deren Anfänge schon auf das Jahr 1805 zurückgehen, ihrem Abschluss näher kommt.
Sie wird noch mindestens zwei umfangreiche Abhandlungen erfordern. Vorerst werde ich es aber bei einer belassen und die erste freie Zeit, die ich danach wieder habe, wieder einem anderen Gegenstand widmen, wahrscheinlich den theoretischen Methoden der höheren Geodäsie.
Es ist seit geraumer Zeit mein Schicksal, immer Arbeiten aufzunehmen, bei denen man in der Darstellung nicht \textls[120]{schnell} vorankommt.\quad\mbox{.\ .\ .\ .\ .\ .}
